% TOP_PROBLEM: {"id": 2399, "title": "Erdős Problem #945", "category_id": 1, "category": {"id": 1, "name": "number_theory", "display_name": "Number Theory", "description": "Properties of integers, prime numbers, Diophantine equations.", "slug": "number-theory", "order_index": 1, "created_at": "2026-07-31T15:26:25.670Z"}, "status": "open"}
% FIRST_POSTED: null
% ENTRY_KIND: statement_only
% Imported from: unsolvedmath_additional_500.tex; UnsolvedMath ID 2399
% !TeX program = lualatex
% Compile twice: lualatex 2399.tex
% Self-contained: no external bibliography, images, or \input files.
% Numeric section labels and visible numbers are original UnsolvedMath IDs.
\documentclass[11pt,a4paper]{article}
\usepackage[margin=24mm,headheight=14pt]{geometry}
\usepackage{fontspec}
\setmainfont{Latin Modern Roman}
\setsansfont{Latin Modern Sans}
\setmonofont{Latin Modern Mono}
\usepackage{amsmath,amssymb,mathtools}
\usepackage{unicode-math}
\setmathfont{Latin Modern Math}
\usepackage{microtype}
\usepackage{enumitem,xurl,needspace}
\usepackage[dvipsnames]{xcolor}
\usepackage{hyperref}
\hypersetup{unicode=true,colorlinks=true,linkcolor=MidnightBlue,urlcolor=MidnightBlue,
 pdftitle={UnsolvedMath Problem 2399},pdfauthor={Source-annotated compilation},bookmarksnumbered=true}
\usepackage{bookmark,tocloft,titlesec,fancyhdr}
\setlength{\cftsecnumwidth}{4.2em}
\cftsetpnumwidth{2.5em}
\cftsetrmarg{4em}
\titleformat{\section}{\Large\bfseries\raggedright}{\thesection}{0.7em}{}
\pagestyle{fancy}\fancyhf{}
\fancyhead[L]{\small\sffamily Additional UnsolvedMath statements}
\fancyhead[R]{\small Original catalogue IDs}
\fancyfoot[C]{\thepage}
\setlength{\parindent}{0pt}
\setlength{\parskip}{5pt plus 1pt}
\setlength{\emergencystretch}{3em}
\setcounter{tocdepth}{1}\setcounter{secnumdepth}{1}
\setlist[itemize]{leftmargin=3em,itemsep=4pt,topsep=4pt,parsep=0pt}
\allowdisplaybreaks
\newcommand{\N}{\mathbb{N}}
\newcommand{\Z}{\mathbb{Z}}
\newcommand{\Q}{\mathbb{Q}}
\newcommand{\R}{\mathbb{R}}
\newcommand{\C}{\mathbb{C}}
\newcommand{\F}{\mathbb{F}}
\newcommand{\A}{\mathbb{A}}
\newcommand{\PP}{\mathbb{P}}
\newcommand{\E}{\mathbb{E}}
\newcommand{\eps}{\varepsilon}
\newcommand{\dd}{\,\mathrm d}
\newcommand{\sref}[2]{\hyperlink{source:#1:#2}{[S#2]}}
\newif\ifshowcatalogue
\showcataloguetrue
\newcommand{\cataloguescope}[1]{\ifshowcatalogue
 \paragraph{Statement recorded in the supplied source.}
 #1\fi}
\begin{document}
% UnsolvedMath ID 2399; source catalogue code EP-945

\setcounter{section}{2398}

\section{Runs with Pairwise Distinct Divisor Counts}\label{2399}

{\small\sffamily UnsolvedMath ID 2399 · Number Theory}

\subsection{Definitions and mathematical statement}

\paragraph{Shared notation.}Unless a section says otherwise, $\N=\{1,2,\ldots\}$,
$\Z$ is the ring of integers, $\Q,\R,\C$ are the rational, real, and complex
fields, $[N]=\{1,\ldots,\lfloor N\rfloor\}$, and $\log$ is the natural
logarithm. For real functions of a parameter tending to infinity,
$f\ll g$ and $f=O(g)$ mean $|f|\leq Cg$ eventually for some fixed $C$;
$f\gg g$ reverses this comparison; $f=o(g)$ means $f/g\to0$;
$f\sim g$ means $f/g\to1$; and $f\asymp g$ means both order bounds.
A subscript on $O$ or $\ll$ specifies allowed constant dependence.
The expression $n^{o(1)}$ in an upper bound means growth smaller than
$n^\epsilon$ for each fixed $\epsilon>0$. Local definitions govern any
exception, finite-field convention, or use of the same symbol for another
object. An asymptotic claim is not automatically an effective algorithm.



The divisor-count function is $\tau(n)=\#\{d\in\N:d\mid n\}$. Divisors are positive, including one and $n$; the function is unrelated to a graph transversal number using the same letter elsewhere. \sref{2399}{1} \sref{2399}{2}

\paragraph{Statement recorded in the supplied source.}
Let $F(x)$ be the maximal $k$ such that there exist $n+1,\ldots,n+k\leq x$ with $\tau(n+1),\ldots,\tau(n+k)$ all distinct (where $\tau(m)$ counts the divisors of $m$). Estimate $F(x)$. In particular, is it true that $ F(x) \leq (\log x)^{O(1)}? $ In other words, is there a constant $C>0$ such that, for all large $x$, every interval $[x,x+(\log x)^C]$ contains two integers with the same number of divisors? \sref{2399}{1}

\paragraph{Residual task reported by the archive.}

The embedded review (2026-08-17) records: Sharpen the upper bound to a fixed log power. \sref{2399}{1}

\subsection{Short English statement}

How long can a consecutive run of integers have all divisor counts different, and is every such run bounded by a fixed logarithmic power of its location?

\subsection{Sources}

{\small

\begin{itemize}

\item[\hypertarget{source:2399:1}{S1}] ULAM.AI, \emph{UnsolvedMath}, supplied \texttt{problems.json}, record 2399 (EP-945), statement and background. The embedded literature-review date is 2026-08-17; that is the archive’s review date, not a fresh certification. Dataset landing page: \url{https://huggingface.co/datasets/ulamai/UnsolvedMath}.

\item[\hypertarget{source:2399:2}{S2}] T. F. Bloom, \emph{Erdős Problems}, Problem 945, \url{https://www.erdosproblems.com/945}. Reference imported from the supplied record; the live problem page was not independently checked for this compilation.

\item[\hypertarget{source:2399:3}{S3}] Erd\H{o}s, P., Some problems and results in number theory. Number theory and combinatorics. Japan 1984 (Tokyo, Okayama and Kyoto, 1984) (1985), 65-87. Bibliographic information imported from the supplied record.

\item[\hypertarget{source:2399:4}{S4}] Erd\H{o}s, P. and Mirsky, L., The distribution of values of the divisor function {$d(n)$}. Proc. London Math. Soc. (3) (1952), 257--271. Bibliographic information imported from the supplied record.

\item[\hypertarget{source:2399:5}{S5}] Guy, Richard K., Unsolved problems in number theory. (2004), xviii+437. Bibliographic information imported from the supplied record.

\end{itemize}

}

\label{end:2399}
\end{document}
